\section{Combien de capital faut-il~?}\label{sec:capital}
% Chapitre 4 (tache 10, passage chapitres 04-05). Figures: fig_kapital_rendite,
% fig_div_vs_4 (voir docs/figures.md). Aucun chiffre tape: macros de
% data/kennzahlen.tex, generees par scripts/rechnung_retraite.py (kapitel_4).
% Controle croise (specification, section 8): le capital requis de la
% strategie de base est \KapitalMaisonSiebzig et celui de la reference
% \KapitalEtfReferenz, les memes macros que le resume reprendra; ce chapitre ne
% les recalcule jamais. Titres d'axes, echelle en millions (y filter) et
% pourcentages (x filter) passes par l'argument optionnel des macros de figure.

Ce chapitre r\'epond \`a la question~: quel capital, en euros de 2026, faut-il avoir r\'euni le jour
du d\'epart pour qu'il te reste \ZielMonat~€ par mois apr\`es retenues, imp\^ot et assurance maladie~?
Le chapitre~\ref{sec:fiscalite} a donn\'e ce qu'il reste de cent euros de dividende selon leur
provenance~; il suffit maintenant de remonter du revenu net au capital. La
section~\ref{sec:capital-methode} pose le calcul, la section~\ref{sec:capital-rendement} montre qu'il
d\'epend d'abord du rendement du dividende, et la section~\ref{sec:capital-regle} compare la strat\'egie
des dividendes seuls avec un retrait de \TauxRetrait~\% par an sur un ETF monde.

\subsection{Du revenu net au capital}\label{sec:capital-methode}

Deux fa\c{c}ons de vivre de son capital sont compar\'ees. Dans la premi\`ere, qui est la th\`ese de ce
plan, tu ne vends jamais rien et tu vis des dividendes vers\'es par ton ETF maison et, selon la
r\'epartition, par un ETF monde distribuant. Dans la seconde, la r\'ef\'erence, tu d\'etiens un ETF
monde capitalisant et tu vends chaque ann\'ee des parts pour \TauxRetrait~\% de la valeur du
portefeuille, la \enquote{r\`egle des \TauxRetrait~\%} que la presse sp\'ecialis\'ee pr\'esente comme
une formule empirique de retrait durable\QP{boerseonline_fire_traum}.

\annahme{Rendement brut du dividende de \RenditeDivMaison~\% pour l'ETF maison (m\'ediane du
portefeuille-exemple du chapitre~\ref{sec:portefeuille}) et de \RenditeDivEtf~\% pour l'ETF monde
distribuant\QH{etf_welt_rendite_div}. Cette seconde valeur est un repli~: le plus grand ETF monde est
capitalisant et ne publie aucun rendement de distribution, et la part distribuante consult\'ee affiche
moins. Si le vrai rendement est plus bas, les r\'epartitions qui gardent une part d'ETF demandent
encore plus de capital que ce qui suit.}

\annahme{R\`egles fiscales et sociales de 2026, appliqu\'ees \`a leur niveau de 2026 et sans dimension
temporelle~: le capital de ce chapitre est celui qu'il faudrait si tu partais aujourd'hui. Imp\^ot au
taux forfaitaire, sans la G\"unstigerpr\"ufung de la section~\ref{sec:fisc-guenstiger}, ce qui rend
ces montants prudents. Les seuils de l'assurance maladie \'etant tenus constants en euros de 2026
(chapitre~\ref{sec:fiscalite}), ce capital vaut aussi pour un d\'epart lointain, \`a l'\'erosion du
forfait d'\'epargnant pr\`es, n\'egligeable \`a ce niveau de revenu. Pour la r\'ef\'erence, tout le retrait est trait\'e comme une plus-value
imposable, cas le plus d\'efavorable, avec l'exon\'eration partielle des ETF sur l'imp\^ot et sur
l'assiette de la cotisation maladie.}

Pour une strat\'egie de dividendes, le capital n\'ecessaire s'obtient par
l'\'equation~\eqref{eq:capital-dividendes-brut}, qui divise le dividende brut annuel dont tu as besoin
par le rendement brut moyen du portefeuille.
\begin{equation}
  K = \frac{B}{m\,r_M + (1-m)\,r_E} \text{,} \label{eq:capital-dividendes-brut}
\end{equation}
o\`u $K$ est le capital en euros de 2026, $B$ le dividende brut annuel qui laisse \ZielJahr~€ apr\`es
retenues, imp\^ot et cotisations, $m$ la part du capital plac\'ee dans l'ETF maison, $r_M$ le rendement
brut de l'ETF maison et $r_E$ celui de l'ETF monde distribuant. Le brut $B$ d\'epend lui-m\^eme de
$m$, puisque l'ETF est moins impos\'e qu'une action d\'etenue en direct~; le mod\`ele r\'esout donc les
deux ensemble. Pour un portefeuille enti\`erement en maison, $B$ est le brut de la cascade du
chapitre~\ref{sec:depart}, \BruttoNoetig~€, ce qui donne \KapitalMaisonHundert~€.

Pour la r\'ef\'erence, l'\'equation~\eqref{eq:capital-retrait-etf} remplace le rendement du dividende
par le taux de retrait.
\begin{equation}
  K = \frac{W}{\tau} \text{,} \label{eq:capital-retrait-etf}
\end{equation}
o\`u $W$ est le retrait brut annuel qui laisse \ZielJahr~€ apr\`es imp\^ot et cotisations, soit
\RetraitBrutEtf~€, et $\tau$ le taux de retrait de \TauxRetrait~\%. Le capital de la r\'ef\'erence vaut
ainsi \KapitalEtfReferenz~€. Les deux \'equations ont la m\^eme forme~; ce qui les s\'epare, c'est que
$\tau$ est nettement plus \'elev\'e que les rendements $r_M$ et $r_E$, et que l'ETF est moins
impos\'e.

\subsection{Le rendement du dividende d\'ecide de tout}\label{sec:capital-rendement}

L'\'equation~\eqref{eq:capital-dividendes-brut} montre que le capital varie comme l'inverse du
rendement. La figure~\ref{fig:capital-selon-rendement} fait varier le rendement brut de l'ETF maison
de \RenditeGrilleBas~\% \`a \RenditeGrilleHaut~\%, le rendement de l'ETF monde restant fix\'e \`a
\RenditeDivEtf~\%.

\linienfigur{fig_kapital_rendite}[xlabel={Rendement brut du dividende de l'ETF maison (\%)},
  ylabel={Capital requis (millions d'euros de 2026)},
  x filter/.expression={x*1e2}, y filter/.expression={y/1e6}, xtick=data,
  legend pos=north east, legend cell align=left, legend style={fill=white}]%
  {rendite}{maison50,maison70,\ColonneToutMaison}%
  {{Moiti\'e en maison},{70~\% en maison},{Tout en maison}}%
  {Capital requis pour \ZielMonat~€ nets par mois de dividendes seuls, selon le rendement brut de
  l'ETF maison, pour trois r\'epartitions entre ETF maison et ETF monde distribuant (r\`egles de 2026,
  euros de 2026).\label{fig:capital-selon-rendement}}

Ce graphique montre que, pour un portefeuille enti\`erement en maison, le capital requis passe de
\KapitalMaisonHundertRenditeBas~€ \`a un rendement de \RenditeGrilleBas~\% \`a
\KapitalMaisonHundertRenditeHaut~€ \`a \RenditeGrilleHaut~\%~: tripler le rendement divise presque le
capital par trois. Au bas de la grille, les trois courbes se rejoignent, car l'ETF monde, moins
impos\'e, compense alors un rendement \`a peine plus faible. Plus le rendement de l'ETF maison
monte, plus la part en maison r\'eduit le capital.

Il faut lire cette figure avec prudence. Un rendement \'elev\'e n'\'economise du capital que s'il dure~:
le chapitre~\ref{sec:risques} montre que des dividendes ont d\'ej\`a baiss\'e fortement pendant des crises, et un
rendement tr\`es sup\'erieur \`a celui du march\'e signale souvent un dividende que le march\'e juge
fragile. Pour que le portefeuille enti\`erement en maison demande le m\^eme capital que la r\'ef\'erence,
il faudrait un rendement brut durable d'environ \RenditeEquilibreEtf~\%, contre \RenditeDivMaison~\%
pour le portefeuille-exemple.

\subsection{Dividendes seuls contre r\`egle des \TauxRetrait~\%}\label{sec:capital-regle}

La figure~\ref{fig:dividendes-contre-retrait} r\'eunit les quatre strat\'egies au rendement de
r\'ef\'erence~: trois r\'epartitions qui ne vivent que des dividendes, et l'ETF monde avec un retrait
de \TauxRetrait~\% par an.

% Ordre des etiquettes = ordre des lignes de data/fig_div_vs_4.csv (MaisonFuenfzig,
% MaisonSiebzig, MaisonHundert, EtfReferenz).
\balkenfigur{fig_div_vs_4}[ylabel={Capital requis (millions d'euros de 2026)},
  y filter/.expression={y/1e6}, ymin=0, enlarge y limits={upper, value=0.2}, bar width=24pt,
  xticklabels={{Moiti\'e en maison},{70~\% en maison},{Tout en maison},{ETF, retrait \TauxRetrait~\%}},
  nodes near coords, every node near coord/.append style={font=\scriptsize,
    /pgf/number format/fixed, /pgf/number format/fixed zerofill, /pgf/number format/precision=2}]%
  {strategie}{kapital}%
  {Capital requis pour \ZielMonat~€ nets par mois~: dividendes seuls selon la r\'epartition, contre
  retrait de \TauxRetrait~\% par an sur un ETF monde capitalisant (r\`egles de 2026, euros de
  2026).\label{fig:dividendes-contre-retrait}}

Ce graphique montre le r\'esultat le moins favorable \`a la th\`ese de ce plan~: vivre des seuls
dividendes demande beaucoup plus de capital qu'un retrait de \TauxRetrait~\%. La r\'epartition du
sc\'enario de base, \ReferenzStrategie{} entre ETF maison et ETF monde, exige \KapitalMaisonSiebzig~€, soit \KapitalDifferenzSiebzigEuro~€ ou
\KapitalDifferenzSiebzigProzent~\% de plus que les \KapitalEtfReferenz~€ de la r\'ef\'erence. M\^eme
la r\'epartition la plus favorable, tout en maison avec \KapitalMaisonHundert~€, reste
\KapitalDifferenzProzent~\% au-dessus, et la moiti\'e en maison monte \`a \KapitalMaisonFuenfzig~€.
Exprim\'e en ann\'ees de d\'epenses, le capital de base repr\'esente \MultipleMaisonSiebzig~fois tes
\ZielJahr~€ annuels, contre \MultipleEtfReferenz~fois pour la r\'ef\'erence.

L'\'ecart a deux causes, d\'ej\`a visibles dans les \'equations~\eqref{eq:capital-dividendes-brut}
et~\eqref{eq:capital-retrait-etf}. La premi\`ere est le taux~: un retrait de \TauxRetrait~\% sort
chaque ann\'ee davantage du portefeuille qu'un dividende de \RenditeDivMaison~\%, parce qu'il puise
aussi dans la hausse des cours. La seconde est l'imp\^ot~: le chapitre~\ref{sec:fiscalite} a montr\'e
qu'un euro retir\'e d'un ETF actions co\^ute moins d'imp\^ot et de cotisation maladie qu'un euro de
dividende d'action d\'etenue en direct.

\bk{Cette comparaison ne dit pas que la r\`egle des \TauxRetrait~\% est sans risque.} Elle vend des
parts chaque ann\'ee, y compris apr\`es une chute des cours, et c'est une r\`egle empirique, pas une
garantie. Finanztip, plus prudent, conseille de retirer chaque ann\'ee d'un \`a trois pour cent du
capital r\'euni au d\'epart en retraite\QP{finanztip_geldanlage}. \`A \TauxRetraitPrudent~\%, la r\'ef\'erence demanderait
\KapitalEtfPrudent~€, et le portefeuille enti\`erement en maison ne la d\'epasserait plus que de
\KapitalDifferenzPrudentProzent~\%. L'argument en faveur des dividendes est que le capital reste
investi et qu'aucune part n'est vendue en pleine crise\QP{justetf_entnahmestrategien}~; l'argument
contraire, que ce confort est surtout psychologique et se paie en diversification et en
rendement\QP{finanznachrichten_irrweg}. La section~\ref{sec:retraite-quatre} met ces deux positions \`a
l'\'epreuve de l'histoire, et le chapitre~\ref{sec:presse} les confronte \`a la presse.

\bk{Retiens de ce chapitre que le capital d\'epend d'abord du rendement du dividende, et que la
strat\'egie des dividendes seuls co\^ute, \`a revenu net \'egal, \KapitalDifferenzSiebzigEuro~€ de
capital de plus que la r\'ef\'erence dans la r\'epartition de base.} Le chapitre~\ref{sec:accumulation}
calcule \`a quel \^age ton \'epargne atteint ce capital, strat\'egie par strat\'egie, et traduit cet
\'ecart en ann\'ees~: \`a \SparquoteZwanzig~\% d'\'epargne, la r\'ef\'erence atteint l'objectif \`a
\AlterEtfReferenzZwanzig~ans alors que la r\'epartition de base ne l'atteint jamais dans l'horizon de
calcul.

\Yannbox{\AnneeYannBasis}{Pour vivre de ses seuls dividendes avec la r\'epartition \ReferenzStrategie{}
entre ETF maison et ETF monde, Yann doit r\'eunir \KapitalMaisonSiebzig~€ en euros de 2026, soit
\MultipleMaisonSiebzig~ans de d\'epenses. Avec un ETF monde et un retrait de \TauxRetrait~\% par an,
\KapitalEtfReferenz~€ suffiraient. En \'epargnant \ReferenzSparquote~\% de son salaire, il franchit
le premier montant en \AnneeYannBasis, avec \KapitalYannBasis~€, soit \KapitalYannEcartCibleProzent~\%
\KapitalYannSensCible{} capital requis~: l'objectif se mesure en fin d'ann\'ee, et la derni\`ere
ann\'ee d'\'epargne fait passer le capital d'un peu moins \`a un peu plus que ce montant.}
